\section{方法总览}
\projectname{} 的核心思想是：不再把谱图看作需要直接翻译成分子字符串的“语言 token”，而是将其视为对分子图、几何与物理响应的多源观测。完整概念路径可写为
\[
\mathcal{S}_{\mathrm{multi}}
\rightarrow
(c, C)
\rightarrow
\mathcal{G}_{\mathrm{cand}}
\rightarrow
D
\rightarrow
R^{(0)}
\rightarrow
R
\rightarrow
\hat{\mathcal{S}}
\rightarrow
\mathrm{Top\mbox{-}K},
\]
其中 $\mathcal{S}_{\mathrm{multi}}$ 为多模态谱图，$(c, C)$ 分别表示全局条件向量与 token 级条件序列，$\mathcal{G}_{\mathrm{cand}}$ 为候选分子集合，$D$ 为 pairwise distance 矩阵，$R^{(0)}$ 为距离几何恢复的初始 3D 构象，$R$ 为等变精修后的几何表示，$\hat{\mathcal{S}}$ 为候选分子的前向预测谱图。

当前 demo 仅显式实现到 $\mathcal{G}_{\mathrm{cand}} \rightarrow \hat{\mathcal{S}} \rightarrow \mathrm{Top\mbox{-}K}$ 的最小闭环；但从文档设计上，后续 2D graph diffusion、pairwise distance 与 TFN-Transformer 的接口已经一并预留。这样既可以在短周期内提交一个可跑通的原型，也能展示项目向更完整科学模型演化的技术路线。

\begin{figure}[htbp]
  \centering
  \placeholderfigure{Figure placeholder: overall architecture}
  \caption{Overall architecture of MultiSpec-GeoDiff: multi-modal spectra are encoded, candidate molecules are generated or retrieved, forward spectra are predicted, and Top-K candidates are reranked.}
\end{figure}

\begin{keybox}[方法分层]
\begin{enumerate}
  \item \textbf{谱图层}：尊重真实横坐标与模态差异；
  \item \textbf{候选层}：先得到可验证的候选集合，而非单点猜测；
  \item \textbf{几何层}：用 pairwise distance 连接 2D 图与 3D 结构；
  \item \textbf{验证层}：通过前向谱图一致性与化学规则完成后验重排。
\end{enumerate}
\end{keybox}
